% GENERATED FILE. Edit canonical lessons, then run npm run generate:books.
% canonical-source-hash: fnv1a64:65c9978468040f5f
% canonical-lessons: JA-C354-doshaburi, JA-C354-kanpa, JA-C354-shikke, JA-C354-higasa, JA-C354-arakajime

\chapter{Weather Plans}
\label{ch:ja-weather-plans}
\hlchaptermodality{354}

\begin{chapteropening}
\textbf{By the end of this chapter:} \emph{I can say what we will do if the weather turns bad.}

Complete the last lesson of chapter 354: I can say what we will do if the weather turns bad.
\end{chapteropening}

\section[\texorpdfstring{\ja{どしゃぶり} (doshaburi)}{どしゃぶり (doshaburi)}]{\ja{どしゃぶり} (doshaburi) — a downpour}

\label{lesson:JA-C354-doshaburi}



\begin{quote}
Before the new one: say the Japanese for a water outage, then the Japanese for tentatively.
\end{quote}

\subsection*{You'll want to know: \ja{どしゃぶり}}
\textbf{\ja{どしゃぶり}} — \emph{doshaburi} — \textquotedblleft{}a downpour\textquotedblright{}.

\begin{quote}
\textbf{\ja{どしゃぶりに} \ja{なったら、しあいは} \ja{ちゅうしでしょう。}} — \emph{doshaburi ni nattara, shiai wa chūshi deshō} — \textquotedblleft{}If it pours, the match will probably be called off.\textquotedblright{}
\end{quote}

\begin{cousinweb}
\textbf{\ja{どしゃぶり}} (\emph{doshaburi}) is built from the Sino-Japanese \emph{dosha}, written with two Chinese characters read with their on'yomi, readings borrowed from Chinese, \emph{do}, \textquotedblleft{}soil\textquotedblright{} (Middle Chinese \emph{thuX}) and \emph{sha}, \textquotedblleft{}sand\textquotedblright{} (\emph{srae}), and the native \emph{buri}, from \emph{furu}, \textquotedblleft{}to fall\textquotedblright{}, voiced from \emph{furi}. JMdict marks that spelling as ateji, characters chosen for their sound, so where \emph{dosha} came from is not settled (KANJIDIC2; JMdict; Middle Chinese in Baxter's transcription).
\end{cousinweb}

\begin{grammarlens}[title={What you've built}]
One more word for saying what happens if.
\end{grammarlens}

\subsection*{Your turn}
\begin{itemize}
\raggedright
  \item \emph{Say it:} \emph{doshaburi}
  \item \emph{Say it:} \emph{doshaburi}, once more
  \item \emph{Say it:} \emph{doshaburi ni nattara, shiai wa chūshi deshō} — If it pours, the match will probably be called off.
  \item \emph{Recall:} say the Japanese for a water outage, then the Japanese for tentatively, then say \emph{doshaburi} again
\end{itemize}

\begin{tcolorbox}[breakable,colback=teal!4,colframe=teal!35!black,title={Before you move on}]
What does \ja{どしゃぶり} mean? (\textquotedblleft{}A downpour\textquotedblright{}.) Say it once more.
\end{tcolorbox}
\section[\texorpdfstring{\ja{かんぱ} (kanpa)}{かんぱ (kanpa)}]{\ja{かんぱ} (kanpa) — a cold wave; a cold snap}

\label{lesson:JA-C354-kanpa}



\begin{quote}
Before the new one: say the Japanese for tentatively, then the Japanese for a downpour.
\end{quote}

\subsection*{You'll want to know: \ja{かんぱ}}
\textbf{\ja{かんぱ}} — \emph{kanpa} — \textquotedblleft{}a cold wave; a cold snap\textquotedblright{}.

\begin{quote}
\textbf{\ja{かんぱが} \ja{きたら、すいどうが} \ja{こおるでしょう。}} — \emph{kanpa ga kitara, suidō ga kōru deshō} — \textquotedblleft{}If a cold snap comes, the water pipes will probably freeze.\textquotedblright{}
\end{quote}

\begin{cousinweb}
\textbf{\ja{かんぱ}} (\emph{kanpa}) is Sino-Japanese (kango): two Chinese characters read with their on'yomi, readings borrowed from Chinese. They are \emph{kan}, \textquotedblleft{}cold\textquotedblright{} (Middle Chinese \emph{han}) and \emph{pa}, \textquotedblleft{}waves\textquotedblright{} (\emph{pa}) (KANJIDIC2; Middle Chinese in Baxter's transcription).
\end{cousinweb}

\begin{grammarlens}[title={What you've built}]
One more word for saying what happens if.
\end{grammarlens}

\subsection*{Your turn}
\begin{itemize}
\raggedright
  \item \emph{Say it:} \emph{kanpa}
  \item \emph{Say it:} \emph{kanpa}, once more
  \item \emph{Say it:} \emph{kanpa ga kitara, suidō ga kōru deshō} — If a cold snap comes, the water pipes will probably freeze.
  \item \emph{Recall:} say the Japanese for tentatively, then the Japanese for a downpour, then say \emph{kanpa} again
\end{itemize}

\begin{tcolorbox}[breakable,colback=teal!4,colframe=teal!35!black,title={Before you move on}]
What does \ja{かんぱ} mean? (\textquotedblleft{}A cold wave; a cold snap\textquotedblright{}.) Say it once more.
\end{tcolorbox}
\section[\texorpdfstring{\ja{しっけ} (shikke)}{しっけ (shikke)}]{\ja{しっけ} (shikke) — dampness; humidity}

\label{lesson:JA-C354-shikke}



\begin{quote}
Before the new one: say the Japanese for a downpour, then the Japanese for a cold snap.
\end{quote}

\subsection*{You'll want to know: \ja{しっけ}}
\textbf{\ja{しっけ}} — \emph{shikke} — \textquotedblleft{}dampness; humidity\textquotedblright{}.

\begin{quote}
\textbf{\ja{つゆに} \ja{なったら、しっけが} \ja{おおく} \ja{なるでしょう。}} — \emph{tsuyu ni nattara, shikke ga ōku naru deshō} — \textquotedblleft{}When the rainy season comes, it will probably get damp.\textquotedblright{}
\end{quote}

\begin{cousinweb}
\textbf{\ja{しっけ}} (\emph{shikke}) is Sino-Japanese (kango): two Chinese characters read with their on'yomi, readings borrowed from Chinese. They are \emph{shitsu}, \textquotedblleft{}damp\textquotedblright{} (Middle Chinese \emph{syip}) and \emph{ke}, \textquotedblleft{}spirit\textquotedblright{} (\emph{khj+jH}) (KANJIDIC2; Middle Chinese in Baxter's transcription). The two join as \emph{shik-ke}.
\end{cousinweb}

\begin{grammarlens}[title={What you've built}]
One more word for saying what happens if.
\end{grammarlens}

\subsection*{Your turn}
\begin{itemize}
\raggedright
  \item \emph{Say it:} \emph{shikke}
  \item \emph{Say it:} \emph{shikke}, once more
  \item \emph{Say it:} \emph{tsuyu ni nattara, shikke ga ōku naru deshō} — When the rainy season comes, it will probably get damp.
  \item \emph{Recall:} say the Japanese for a downpour, then the Japanese for a cold snap, then say \emph{shikke} again
\end{itemize}

\begin{tcolorbox}[breakable,colback=teal!4,colframe=teal!35!black,title={Before you move on}]
What does \ja{しっけ} mean? (\textquotedblleft{}Dampness; humidity\textquotedblright{}.) Say it once more.
\end{tcolorbox}
\section[\texorpdfstring{\ja{ひがさ} (higasa)}{ひがさ (higasa)}]{\ja{ひがさ} (higasa) — a parasol}

\label{lesson:JA-C354-higasa}



\begin{quote}
Before the new one: say the Japanese for a cold snap, then the Japanese for dampness.
\end{quote}

\subsection*{You'll want to know: \ja{ひがさ}}
\textbf{\ja{ひがさ}} — \emph{higasa} — \textquotedblleft{}a parasol\textquotedblright{}.

\begin{quote}
\textbf{\ja{はれたら、ひがさを} \ja{もって} \ja{ください。}} — \emph{haretara, higasa o motte kudasai} — \textquotedblleft{}If it is sunny, please take a parasol.\textquotedblright{}
\end{quote}

\begin{cousinweb}
\textbf{\ja{ひがさ}} (\emph{higasa}) is a native Japanese word (wago), inherited rather than borrowed from Chinese: it is built from \emph{hi}, \textquotedblleft{}sun\textquotedblright{} and \emph{gasa}, \textquotedblleft{}umbrella (kasa)\textquotedblright{}, and its kanji were attached later as kun'yomi (JMdict).
\end{cousinweb}

\begin{grammarlens}[title={What you've built}]
One more word for saying what happens if.
\end{grammarlens}

\subsection*{Your turn}
\begin{itemize}
\raggedright
  \item \emph{Say it:} \emph{higasa}
  \item \emph{Say it:} \emph{higasa}, once more
  \item \emph{Say it:} \emph{haretara, higasa o motte kudasai} — If it is sunny, please take a parasol.
  \item \emph{Recall:} say the Japanese for a cold snap, then the Japanese for dampness, then say \emph{higasa} again
\end{itemize}

\begin{tcolorbox}[breakable,colback=teal!4,colframe=teal!35!black,title={Before you move on}]
What does \ja{ひがさ} mean? (\textquotedblleft{}A parasol\textquotedblright{}.) Say it once more.
\end{tcolorbox}
\section[\texorpdfstring{\ja{あらかじめ} (arakajime)}{あらかじめ (arakajime)}]{\ja{あらかじめ} (arakajime) — in advance; beforehand}

\label{lesson:JA-C354-arakajime}



\begin{quote}
Before the new one: say the Japanese for dampness, then the Japanese for a parasol.
\end{quote}

\subsection*{You'll want to know: \ja{あらかじめ}}
\textbf{\ja{あらかじめ}} — \emph{arakajime} — \textquotedblleft{}in advance; beforehand\textquotedblright{}.

\begin{quote}
\textbf{\ja{よていが} \ja{かわったら、あらかじめ} \ja{おしえて} \ja{ください。}} — \emph{yotei ga kawattara, arakajime oshiete kudasai} — \textquotedblleft{}If the plans change, please let me know in advance.\textquotedblright{}
\end{quote}

\begin{cousinweb}
\textbf{\ja{あらかじめ}} (\emph{arakajime}) is a native Japanese word (wago), inherited rather than borrowed from Chinese, and its older origin is not settled. JMdict notes that it is usually written in kana (JMdict).
\end{cousinweb}

\begin{grammarlens}[title={What you've built}]
One more word for saying what happens if.
\end{grammarlens}

\subsection*{Your turn}
\begin{itemize}
\raggedright
  \item \emph{Say it:} \emph{arakajime}
  \item \emph{Say it:} \emph{arakajime}, once more
  \item \emph{Say it:} \emph{yotei ga kawattara, arakajime oshiete kudasai} — If the plans change, please let me know in advance.
  \item \emph{Recall:} say the Japanese for dampness, then the Japanese for a parasol, then say \emph{arakajime} again
\end{itemize}

\begin{tcolorbox}[breakable,colback=teal!4,colframe=teal!35!black,title={Before you move on}]
What does \ja{あらかじめ} mean? (\textquotedblleft{}In advance; beforehand\textquotedblright{}.) Say it once more.
\end{tcolorbox}
